% Options for packages loaded elsewhere
\PassOptionsToPackage{unicode}{hyperref}
\PassOptionsToPackage{hyphens}{url}
\PassOptionsToPackage{dvipsnames,svgnames,x11names}{xcolor}
\documentclass[
  conference]{IEEEtran}
\usepackage{xcolor}
\usepackage{amsmath,amssymb}
\setcounter{secnumdepth}{-\maxdimen} % remove section numbering
\usepackage{iftex}
\ifPDFTeX
  \usepackage[T1]{fontenc}
  \usepackage[utf8]{inputenc}
  \usepackage{textcomp} % provide euro and other symbols
\else % if luatex or xetex
  \usepackage{unicode-math} % this also loads fontspec
  \defaultfontfeatures{Scale=MatchLowercase}
  \defaultfontfeatures[\rmfamily]{Ligatures=TeX,Scale=1}
\fi
\usepackage{lmodern}
\ifPDFTeX\else
  % xetex/luatex font selection
\fi
% Use upquote if available, for straight quotes in verbatim environments
\IfFileExists{upquote.sty}{\usepackage{upquote}}{}
\IfFileExists{microtype.sty}{% use microtype if available
  \usepackage[]{microtype}
  \UseMicrotypeSet[protrusion]{basicmath} % disable protrusion for tt fonts
}{}
\makeatletter
\@ifundefined{KOMAClassName}{% if non-KOMA class
  \IfFileExists{parskip.sty}{%
    \usepackage{parskip}
  }{% else
    \setlength{\parindent}{0pt}
    \setlength{\parskip}{6pt plus 2pt minus 1pt}}
}{% if KOMA class
  \KOMAoptions{parskip=half}}
\makeatother
\usepackage{longtable,booktabs,array}
\usepackage{caption}
\captionsetup[table]{skip=6pt}
\newcounter{none} % for unnumbered tables
\usepackage{calc} % for calculating minipage widths
% Correct order of tables after \paragraph or \subparagraph
\usepackage{etoolbox}
\makeatletter
\patchcmd\longtable{\par}{\if@noskipsec\mbox{}\fi\par}{}{}
\makeatother
% definitions for citeproc citations
\NewDocumentCommand\citeproctext{}{}
\NewDocumentCommand\citeproc{mm}{%
  \begingroup\def\citeproctext{#2}\cite{#1}\endgroup}
\makeatletter
 % allow citations to break across lines
 \let\@cite@ofmt\@firstofone
 % avoid brackets around text for \cite:
 \def\@biblabel#1{}
 \def\@cite#1#2{{#1\if@tempswa , #2\fi}}
\makeatother
\newlength{\cslhangindent}
\setlength{\cslhangindent}{1.5em}
\newlength{\csllabelwidth}
\setlength{\csllabelwidth}{3em}
\newenvironment{CSLReferences}[2] % #1 hanging-indent, #2 entry-spacing
 {\begin{list}{}{%
  \setlength{\itemindent}{0pt}
  \setlength{\leftmargin}{0pt}
  \setlength{\parsep}{0pt}
  % turn on hanging indent if param 1 is 1
  \ifodd #1
   \setlength{\leftmargin}{\cslhangindent}
   \setlength{\itemindent}{-1\cslhangindent}
  \fi
  % set entry spacing
  \setlength{\itemsep}{#2\baselineskip}}}
 {\end{list}}
\usepackage{calc}
\newcommand{\CSLBlock}[1]{\hfill\break\parbox[t]{\linewidth}{\strut\ignorespaces#1\strut}}
\newcommand{\CSLLeftMargin}[1]{\parbox[t]{\csllabelwidth}{\strut#1\strut}}
\newcommand{\CSLRightInline}[1]{\parbox[t]{\linewidth - \csllabelwidth}{\strut#1\strut}}
\newcommand{\CSLIndent}[1]{\hspace{\cslhangindent}#1}
\setlength{\emergencystretch}{3em} % prevent overfull lines
\providecommand{\tightlist}{%
  \setlength{\itemsep}{0pt}\setlength{\parskip}{0pt}}
\usepackage{times}
\usepackage{bookmark}
\IfFileExists{xurl.sty}{\usepackage{xurl}}{} % add URL line breaks if available
\urlstyle{same}
% fallback for those not using the hyperref driver hyperxmp:
\makeatletter
\@ifundefined{xmpquote}{\newcommand{\xmpquote}[1]{#1}}{}
\makeatother
\hypersetup{
  pdftitle={Stability-Penalized Split Selection for Decision Tree Classification: A Reproducible Multi-Dataset Evaluation},
  pdfauthor={Rofiq Fauzi (p36202600047@mhs.dinus.ac.id), Magister Artificial Intelligence, Universitas Dian Nuswantoro},
  colorlinks=true,
  linkcolor={black},
  filecolor={Maroon},
  citecolor={Blue},
  urlcolor={black},
  pdfcreator={LaTeX via pandoc}}

\title{Stability-Penalized Split Selection for Decision Tree
Classification: A Reproducible Multi-Dataset Evaluation}
\author{Rofiq Fauzi (p36202600047@mhs.dinus.ac.id), Magister Artificial
Intelligence, Universitas Dian Nuswantoro}
\date{}

\begin{document}
\maketitle

\textbf{Keywords}--decision trees; split criterion; bootstrap stability;
gain variance; nested cross-validation; nonparametric statistical
comparison.

\subsection{Abstract}\label{abstract}

The CART decision tree is fast and interpretable, but its structure is
unstable: small changes in the training data can alter the selected
split and, ultimately, the predictions. This study evaluates a very
simple candidate formulation, named \emph{Stability-Penalized Decision
Tree} (SPDT), in which each candidate split is scored as Gini gain minus
a weighted penalty on the variance of that gain over bootstrap resamples
of the node data:
\texttt{Score(s)\ =\ Gain(s,D)\ -\ lambda*Var\_b{[}Gain(s,D\^{}(b)){]}}.
A literature search up to 22 October 2026 found no formula-level
equivalent of this criterion, but the tree-stability research space is
crowded; the formulation is therefore positioned as a candidate, not as
a novelty claim. The evaluation uses 8 UCI tabular classification
datasets with repeated nested cross-validation (5 folds x 3 repeats on
three seeds), hyperparameter tuning in the inner loop, and Macro-F1 as
the primary metric. Results: SPDT versus custom CART (lambda=0) records
6 wins, 0 ties, and 2 losses with a mean Macro-F1 difference of +0.0058.
The Friedman test gives chi\^{}2(2)=4.0000, p=0.1353, Kendall's W=0.2500
-- insufficient evidence to reject the null hypothesis, so confirmatory
post-hoc comparisons were not run. Ablations show that the penalty
strength lambda and the number of resamples B barely affect Macro-F1,
tree-depth constraints matter far more, and the prediction-stability
measure of SPDT is slightly worse than that of custom CART at roughly
40x the fitting time. These findings are an honest null result: across
the eight small-to-medium datasets tested, the bootstrap gain-variance
penalty provides no benefit sufficient to outweigh the simplicity of
CART, without ruling out benefits in other data regimes.

\subsection{I. Introduction}\label{i.-introduction}

Decision trees are nonparametric classifiers that recursively partition
the feature space. The CART split criterion selects the candidate
(feature, threshold) that maximizes the reduction of Gini impurity
(Breiman et al. 1984; Quinlan 1986). However, tree structure is known to
be unstable: small perturbations of the training data can change the
selected split (Li and Clarke 2002), and consequently the predictions --
the very phenomenon behind \emph{bagging} and \emph{random forests}
(Breiman 1996, 2001; Dietterich 2000).

Various approaches have reduced this instability at different levels:
significance filtering, structural restriction, ensemble aggregation, or
stable model selection (Briand et al. 2009; Lee et al. 2024; Mirzamomen
and Kangavari 2017; Wang et al. 2017). What is rarely tested directly is
the narrower question:

\begin{quote}
\textbf{Can a bootstrap variance penalty on the split gain, subtracted
directly from the split score, improve out-of-sample Macro-F1 and/or
prediction stability compared with standard CART under repeated
resampling, at acceptable runtime cost?}
\end{quote}

Null hypothesis: no systematic improvement over the baseline across
independent datasets. No hypothesis is made that SPDT \emph{must} win.

The claimed contributions are limited to: (i) an explicit candidate
formulation; (ii) a reproducible implementation and audit; (iii) a
leakage-free nested-CV evaluation; (iv) honest ablation and statistical
analysis of both positive and negative outcomes.

\subsection{II. Related Work}\label{ii.-related-work}

\textbf{Decision tree stability.} Structural tree instability was
formalized by Li and Clarke (Li and Clarke 2002); quantitative stability
and tree-similarity measures were developed, among others, by Turney
(Turney 1995) and Briand et al. (Briand et al. 2009), who also evaluated
bootstrap-based node-level stabilization (the NLS/Dannegger lineage).
Lee et al. (Lee et al. 2024) proposed ANLS, which runs bootstrap
iterations at each node and selects the most frequently chosen predictor
-- the closest resampling-per-node mechanism to this study, but its
aggregation is frequency voting rather than a gain-variance penalty.
Mirzamomen and Kangavari (Mirzamomen and Kangavari 2017) stabilize trees
through multi-attribute split tests.

\textbf{Alternative split criteria.} Two-term Tsallis entropy criteria
(Wang et al. 2017), label-noise-robust criteria (Yang et al. 2019),
bias-corrected Gini split selection (Strobl et al. 2007), and
Hoeffding-bound data-stream trees (Domingos and Hulten 2000) show that
the ``modify the split score'' research space is active. Delise (Delise
2026) scores split consistency across environments in boosting.
Penalizing a resampled score by its variance is also known outside trees
(the mean-variance pattern in model selection). None of the works we
were able to verify contains the formula
\texttt{Gain\ -\ lambda*Var\_bootstrap{[}Gain{]}}; nevertheless, a
negative search result is not proof of novelty, and the formulation
remains a candidate.

\subsection{III. Method}\label{iii.-method}

\subsubsection{A. Proposed split
criterion}\label{a.-proposed-split-criterion}

For a node with dataset D (training portion only), with
\texttt{Gini(D)\ =\ 1\ -\ Sum\_c\ p(c\textbar{}D)\^{}2} and split gain

\texttt{Gain(s,D)\ =\ Gini(D)\ -\ (\textbar{}D\_L\textbar{}/\textbar{}D\textbar{})*Gini(D\_L)\ -\ (\textbar{}D\_R\textbar{}/\textbar{}D\textbar{})*Gini(D\_R)},

B bootstrap resamples \texttt{D\^{}(b)} are drawn from the node rows.
For every valid candidate s:

\texttt{V(s)\ =\ Var\_b{[}Gain(s,D\^{}(b)){]}} (sample variance, ddof=1)

\textbf{Main score: \texttt{Score(s)\ =\ Gain(s,D)\ -\ lambda*V(s)}}.

Deterministic policies: candidate thresholds are midpoints between
sorted distinct adjacent values;
\texttt{min\_samples\_leaf}/\texttt{min\_samples\_split} are applied to
the original node counts; an empty child partition inside a resample
contributes 0; leaves predict the majority class with smallest-label
tie-breaking; bootstrap seeds derive from
\texttt{(random\_state,\ node\ id)}. For lambda=0 the entire bootstrap
computation is skipped and the estimator degenerates to the custom CART
reference (verified identical to an independent reference implementation
in unit tests). A separate exploratory variant (normalization by
\texttt{V/(Gain\^{}2+eps)}) is documented but is not the main method. An
important limitation: V(s) is scale-dependent on the gain.

\subsubsection{B. Evaluation design}\label{b.-evaluation-design}

Datasets: 8 UCI tabular classification datasets (Table I). Protocol:
repeated stratified nested CV -- outer 5 folds x 3 repeats for each of
the seeds {[}42, 123, 2026{]} (45 evaluations per dataset-method,
identical splits for all methods); inner stratified 3-fold tuning with
deterministically derived seeds. Preprocessing is median imputation
fitted on the training portion only (pipeline). Primary metric Macro-F1;
secondary: balanced accuracy, accuracy, fit/predict time, node count,
tree depth. Prediction stability is measured separately as mean pairwise
disagreement among models trained on 10 subsamples (90\%) of each outer
training fold, evaluated on the same outer test fold -- reported
separately from accuracy.

\begin{longtable}[]{@{}lllll@{}}
\caption{Dataset inventory.}\tabularnewline
\toprule\noalign{}
Dataset & N & Features & Classes & Source \\
\midrule\noalign{}
\endfirsthead
\toprule\noalign{}
Dataset & N & Features & Classes & Source \\
\midrule\noalign{}
\endhead
\bottomrule\noalign{}
\endlastfoot
banknote & 1372 & 4 & 2 & ucimlrepo \\
glass & 214 & 9 & 6 & ucimlrepo \\
ionosphere & 351 & 34 & 2 & ucimlrepo \\
iris & 150 & 4 & 3 & ucimlrepo \\
seeds & 210 & 7 & 3 & uci\_static\_zip \\
sonar & 208 & 60 & 2 & ucimlrepo \\
wdbc & 569 & 30 & 2 & ucimlrepo \\
wine & 178 & 13 & 3 & ucimlrepo \\
\end{longtable}

\subsubsection{C. Statistical analysis}\label{c.-statistical-analysis}

The unit of inference is the dataset (averaged over 45 outer folds), not
the fold. The three primary methods (sklearn CART, custom CART lambda=0,
SPDT) are compared with the Friedman test using datasets as blocks; only
when the omnibus test provides grounds for it are paired Wilcoxon
signed-rank tests with Holm correction and rank-biserial effect sizes
run. p \textgreater= 0.05 is never interpreted as evidence of
equivalence. Random forest is included as a non-equivalent contextual
benchmark and is not part of the primary omnibus test.

\subsection{IV. Results and
Discussion}\label{iv.-results-and-discussion}

\begin{longtable}[]{@{}
  >{\raggedright\arraybackslash}p{(\linewidth - 10\tabcolsep) * \real{0.1667}}
  >{\raggedright\arraybackslash}p{(\linewidth - 10\tabcolsep) * \real{0.1667}}
  >{\raggedright\arraybackslash}p{(\linewidth - 10\tabcolsep) * \real{0.1667}}
  >{\raggedright\arraybackslash}p{(\linewidth - 10\tabcolsep) * \real{0.1667}}
  >{\raggedright\arraybackslash}p{(\linewidth - 10\tabcolsep) * \real{0.1667}}
  >{\raggedright\arraybackslash}p{(\linewidth - 10\tabcolsep) * \real{0.1667}}@{}}
\caption{Macro-F1 mean +/- sd (descriptive, fold-level) per dataset and
method.}\tabularnewline
\toprule\noalign{}
\begin{minipage}[b]{\linewidth}\raggedright
Dataset
\end{minipage} & \begin{minipage}[b]{\linewidth}\raggedright
cart\_sklearn
\end{minipage} & \begin{minipage}[b]{\linewidth}\raggedright
cart\_custom
\end{minipage} & \begin{minipage}[b]{\linewidth}\raggedright
spdt
\end{minipage} & \begin{minipage}[b]{\linewidth}\raggedright
random\_forest
\end{minipage} & \begin{minipage}[b]{\linewidth}\raggedright
Delta (spdt-cart\_custom)
\end{minipage} \\
\midrule\noalign{}
\endfirsthead
\toprule\noalign{}
\begin{minipage}[b]{\linewidth}\raggedright
Dataset
\end{minipage} & \begin{minipage}[b]{\linewidth}\raggedright
cart\_sklearn
\end{minipage} & \begin{minipage}[b]{\linewidth}\raggedright
cart\_custom
\end{minipage} & \begin{minipage}[b]{\linewidth}\raggedright
spdt
\end{minipage} & \begin{minipage}[b]{\linewidth}\raggedright
random\_forest
\end{minipage} & \begin{minipage}[b]{\linewidth}\raggedright
Delta (spdt-cart\_custom)
\end{minipage} \\
\midrule\noalign{}
\endhead
\bottomrule\noalign{}
\endlastfoot
banknote & 0.982 +/- 0.008 & 0.981 +/- 0.007 & 0.981 +/- 0.007 & 0.994
+/- 0.004 & +0.0004 \\
glass & 0.620 +/- 0.097 & 0.609 +/- 0.093 & 0.626 +/- 0.093 & 0.742 +/-
0.075 & +0.0166 \\
ionosphere & 0.873 +/- 0.034 & 0.881 +/- 0.038 & 0.879 +/- 0.038 & 0.929
+/- 0.037 & -0.0027 \\
iris & 0.949 +/- 0.033 & 0.945 +/- 0.038 & 0.944 +/- 0.037 & 0.955 +/-
0.034 & -0.0008 \\
seeds & 0.897 +/- 0.032 & 0.906 +/- 0.036 & 0.913 +/- 0.034 & 0.923 +/-
0.038 & +0.0066 \\
sonar & 0.719 +/- 0.065 & 0.706 +/- 0.073 & 0.725 +/- 0.061 & 0.833 +/-
0.065 & +0.0189 \\
wdbc & 0.925 +/- 0.031 & 0.925 +/- 0.030 & 0.928 +/- 0.030 & 0.956 +/-
0.023 & +0.0037 \\
wine & 0.905 +/- 0.047 & 0.905 +/- 0.044 & 0.909 +/- 0.046 & 0.981 +/-
0.024 & +0.0040 \\
\end{longtable}

\begin{longtable}[]{@{}ll@{}}
\caption{Mean Macro-F1 rank.}\tabularnewline
\toprule\noalign{}
Method & Mean rank \\
\midrule\noalign{}
\endfirsthead
\toprule\noalign{}
Method & Mean rank \\
\midrule\noalign{}
\endhead
\bottomrule\noalign{}
\endlastfoot
spdt & 1.5000 \\
cart\_sklearn & 2.0000 \\
cart\_custom & 2.5000 \\
\end{longtable}

\begin{longtable}[]{@{}
  >{\raggedright\arraybackslash}p{(\linewidth - 12\tabcolsep) * \real{0.1429}}
  >{\raggedright\arraybackslash}p{(\linewidth - 12\tabcolsep) * \real{0.1429}}
  >{\raggedright\arraybackslash}p{(\linewidth - 12\tabcolsep) * \real{0.1429}}
  >{\raggedright\arraybackslash}p{(\linewidth - 12\tabcolsep) * \real{0.1429}}
  >{\raggedright\arraybackslash}p{(\linewidth - 12\tabcolsep) * \real{0.1429}}
  >{\raggedright\arraybackslash}p{(\linewidth - 12\tabcolsep) * \real{0.1429}}
  >{\raggedright\arraybackslash}p{(\linewidth - 12\tabcolsep) * \real{0.1429}}@{}}
\caption{Dataset-level statistical tests.}\tabularnewline
\toprule\noalign{}
\begin{minipage}[b]{\linewidth}\raggedright
Test
\end{minipage} & \begin{minipage}[b]{\linewidth}\raggedright
Statistic
\end{minipage} & \begin{minipage}[b]{\linewidth}\raggedright
df
\end{minipage} & \begin{minipage}[b]{\linewidth}\raggedright
p
\end{minipage} & \begin{minipage}[b]{\linewidth}\raggedright
\#datasets
\end{minipage} & \begin{minipage}[b]{\linewidth}\raggedright
Kendall W
\end{minipage} & \begin{minipage}[b]{\linewidth}\raggedright
Decision
\end{minipage} \\
\midrule\noalign{}
\endfirsthead
\toprule\noalign{}
\begin{minipage}[b]{\linewidth}\raggedright
Test
\end{minipage} & \begin{minipage}[b]{\linewidth}\raggedright
Statistic
\end{minipage} & \begin{minipage}[b]{\linewidth}\raggedright
df
\end{minipage} & \begin{minipage}[b]{\linewidth}\raggedright
p
\end{minipage} & \begin{minipage}[b]{\linewidth}\raggedright
\#datasets
\end{minipage} & \begin{minipage}[b]{\linewidth}\raggedright
Kendall W
\end{minipage} & \begin{minipage}[b]{\linewidth}\raggedright
Decision
\end{minipage} \\
\midrule\noalign{}
\endhead
\bottomrule\noalign{}
\endlastfoot
friedman & 4.0000 & 2 & 0.1353 & 8 & 0.2500 & no omnibus evidence;
post-hoc not run \\
\end{longtable}

\emph{Post-hoc not run: the Friedman omnibus test gave no grounds (p
\textgreater= 0.05), per protocol.}

Table: Pairwise (post-hoc) comparisons.

Paired summary SPDT vs custom CART: \textbf{6 wins / 0 ties / 2 losses},
mean Macro-F1 difference +0.0058.

\begin{longtable}[]{@{}
  >{\raggedright\arraybackslash}p{(\linewidth - 10\tabcolsep) * \real{0.1667}}
  >{\raggedright\arraybackslash}p{(\linewidth - 10\tabcolsep) * \real{0.1667}}
  >{\raggedright\arraybackslash}p{(\linewidth - 10\tabcolsep) * \real{0.1667}}
  >{\raggedright\arraybackslash}p{(\linewidth - 10\tabcolsep) * \real{0.1667}}
  >{\raggedright\arraybackslash}p{(\linewidth - 10\tabcolsep) * \real{0.1667}}
  >{\raggedright\arraybackslash}p{(\linewidth - 10\tabcolsep) * \real{0.1667}}@{}}
\caption{Computational cost.}\tabularnewline
\toprule\noalign{}
\begin{minipage}[b]{\linewidth}\raggedright
Method
\end{minipage} & \begin{minipage}[b]{\linewidth}\raggedright
Fit mean (s)
\end{minipage} & \begin{minipage}[b]{\linewidth}\raggedright
Fit median (s)
\end{minipage} & \begin{minipage}[b]{\linewidth}\raggedright
Predict mean (s)
\end{minipage} & \begin{minipage}[b]{\linewidth}\raggedright
Mean nodes
\end{minipage} & \begin{minipage}[b]{\linewidth}\raggedright
Mean depth
\end{minipage} \\
\midrule\noalign{}
\endfirsthead
\toprule\noalign{}
\begin{minipage}[b]{\linewidth}\raggedright
Method
\end{minipage} & \begin{minipage}[b]{\linewidth}\raggedright
Fit mean (s)
\end{minipage} & \begin{minipage}[b]{\linewidth}\raggedright
Fit median (s)
\end{minipage} & \begin{minipage}[b]{\linewidth}\raggedright
Predict mean (s)
\end{minipage} & \begin{minipage}[b]{\linewidth}\raggedright
Mean nodes
\end{minipage} & \begin{minipage}[b]{\linewidth}\raggedright
Mean depth
\end{minipage} \\
\midrule\noalign{}
\endhead
\bottomrule\noalign{}
\endlastfoot
cart\_custom & 0.5340 & 0.4163 & 0.0009 & 26.6722 & 5.1000 \\
cart\_sklearn & 0.2693 & 0.2605 & 0.0010 & 27.4222 & 5.2917 \\
random\_forest & 1.1169 & 1.0863 & 0.0573 & nan & nan \\
spdt & 21.4302 & 14.6491 & 0.0010 & 28.5944 & 5.4611 \\
\end{longtable}

\begin{longtable}[]{@{}lll@{}}
\caption{Prediction stability (disagreement; not an accuracy
metric).}\tabularnewline
\toprule\noalign{}
Method & Mean disagreement & SD \\
\midrule\noalign{}
\endfirsthead
\toprule\noalign{}
Method & Mean disagreement & SD \\
\midrule\noalign{}
\endhead
\bottomrule\noalign{}
\endlastfoot
cart\_custom & 0.0928 & 0.0853 \\
cart\_sklearn & 0.0983 & 0.0885 \\
random\_forest & 0.0271 & 0.0337 \\
spdt & 0.1014 & 0.0934 \\
\end{longtable}

\begin{longtable}[]{@{}lllll@{}}
\caption{Ablation on identical splits.}\tabularnewline
\toprule\noalign{}
Variant & Mean macro F1 & SD & Fit mean (s) & Mean nodes \\
\midrule\noalign{}
\endfirsthead
\toprule\noalign{}
Variant & Mean macro F1 & SD & Fit mean (s) & Mean nodes \\
\midrule\noalign{}
\endhead
\bottomrule\noalign{}
\endlastfoot
A0\_lam0 & 0.8599 & 0.1270 & 0.0201 & 37.6611 \\
A1\_lam0.1\_B10 & 0.8598 & 0.1256 & 0.0860 & 37.6111 \\
A1\_lam1\_B10 & 0.8601 & 0.1255 & 0.0862 & 37.5278 \\
A2\_B20 & 0.8601 & 0.1251 & 0.1136 & 37.4556 \\
A2\_B5 & 0.8613 & 0.1243 & 0.0711 & 37.4667 \\
A3\_d3\_leaf1 & 0.8357 & 0.1592 & 0.0371 & 12.8722 \\
A3\_d3\_leaf5 & 0.8353 & 0.1546 & 0.0373 & 11.8611 \\
A3\_d8\_leaf1 & 0.8592 & 0.1272 & 0.0821 & 35.2500 \\
A3\_d8\_leaf5 & 0.8521 & 0.1381 & 0.0820 & 24.7389 \\
\end{longtable}

\begin{longtable}[]{@{}llll@{}}
\caption{SPDT hyperparameter selection frequency (45 outer folds per
dataset).}\tabularnewline
\toprule\noalign{}
parameter & value & frequency & share \\
\midrule\noalign{}
\endfirsthead
\toprule\noalign{}
parameter & value & frequency & share \\
\midrule\noalign{}
\endhead
\bottomrule\noalign{}
\endlastfoot
max\_depth & 5.0000 & 140 & 0.3889 \\
max\_depth & 8.0000 & 102 & 0.2833 \\
max\_depth & 3.0000 & 92 & 0.2556 \\
min\_samples\_leaf & 1.0000 & 201 & 0.5583 \\
min\_samples\_leaf & 2.0000 & 97 & 0.2694 \\
min\_samples\_leaf & 5.0000 & 62 & 0.1722 \\
min\_samples\_split & 2.0000 & 252 & 0.7000 \\
min\_samples\_split & 5.0000 & 69 & 0.1917 \\
min\_samples\_split & 10.0000 & 39 & 0.1083 \\
n\_bootstrap & 5.0000 & 191 & 0.5306 \\
n\_bootstrap & 20.0000 & 87 & 0.2417 \\
n\_bootstrap & 10.0000 & 82 & 0.2278 \\
penalty\_lambda & 10.0000 & 233 & 0.6472 \\
penalty\_lambda & 0.0000 & 61 & 0.1694 \\
penalty\_lambda & 0.0100 & 41 & 0.1139 \\
penalty\_lambda & 1.0000 & 22 & 0.0611 \\
penalty\_lambda & 0.1000 & 3 & 0.0083 \\
\end{longtable}

\textbf{Discussion.} The results descriptively favor SPDT but do not
reach confirmatory evidence. SPDT achieves the best mean rank (1.500 vs
2.000 for sklearn CART and 2.500 for custom CART) and wins on 6 of 8
datasets paired against custom CART, yet the Friedman test finds
insufficient evidence of a difference (p = 0.1353), and the exploratory
SPDT-custom CART comparison gives a raw p of 0.078 -- it does not
survive the 5\% level and must not be read as significant. With only 8
datasets, statistical power is limited; per protocol, the results are
reported as they stand and p \textgreater= 0.05 is not interpreted as
equivalence.

A second important finding: the stability penalty does not reduce
prediction disagreement -- SPDT is in fact slightly higher (0.101) than
custom CART (0.093) and sklearn CART (0.098). Random forest remains the
most stable (0.027) through ensemble aggregation. The motivating promise
of the method (better prediction stability) therefore does not
materialize on this benchmark; a gain-variance penalty is no substitute
for ensembling.

A third finding from the identical-split ablations: neither the penalty
strength lambda nor the number of resamples B materially changes
Macro-F1 (all lambda and B variants fall in the 0.859-0.861 range),
whereas depth/leaf constraints matter far more (depth 3 drops to about
0.836). Inner tuning tends to select large lambda (10 in 64.7\% of
folds) and a small B (5 in 53.1\%), indicating the penalty is often
selected but has a small effect on accuracy. The cost is real: SPDT
fitting time is roughly 40x that of custom CART. Overall, the
effectiveness of the bootstrap gain-variance penalty is
dataset-dependent and small in magnitude; simple impurity-based
splitting remains preferable for these small-to-medium datasets.

\subsection{V. Conclusion}\label{v.-conclusion}

This study evaluated the SPDT candidate formulation reproducibly, with
leakage-free repeated nested CV, controlled ablation on identical
splits, and dataset-level statistical analysis on 8 UCI datasets.
Descriptively, SPDT achieved the best mean rank and won 6 of 8 paired
comparisons against custom CART, but the Friedman test provided
insufficient evidence of a difference (chi\^{}2(2)=4.0000, p=0.1353), so
there is no confirmatory basis for claiming improvement. Moreover, the
prediction-stability measure did not improve and was slightly worse,
while compute cost rose roughly 40x. Under the tested conditions, the
simplicity of CART remains preferable. Limitations: eight small
datasets, Gini impurity only, a single disagreement-based stability
proxy, a formula-level novelty claim that is not established, and the
absence of the closest literature criteria (ANLS, minimax split) as
comparators. The most valuable future work: direct comparison with
node-level bootstrap stabilization from the literature, a theoretically
grounded penalty normalization by gain scale, and evaluation on larger
benchmarks.

\subsection{References}\label{references}

\protect\phantomsection\label{refs}
\begin{CSLReferences}{1}{1}
\bibitem[\citeproctext]{ref-breiman1996bagging}
Breiman, Leo. 1996. {``Bagging Predictors.''} \emph{Machine Learning} 24
(2): 123--40. \url{https://doi.org/10.1023/A:1018054314350}.

\bibitem[\citeproctext]{ref-breiman2001randomforests}
Breiman, Leo. 2001. {``Random Forests.''} \emph{Machine Learning} 45
(1): 5--32. \url{https://doi.org/10.1023/A:1010933404324}.

\bibitem[\citeproctext]{ref-breiman1984cart}
Breiman, Leo, Jerome H. Friedman, Richard A. Olshen, and Charles J.
Stone. 1984. \emph{Classification and Regression Trees}. Wadsworth;
expanded edition Chapman \& Hall/CRC, 1993.

\bibitem[\citeproctext]{ref-briand2009stability}
Briand, Bénédicte, Gilles R. Ducharme, Vanessa Parache, and Catherine
Mercat-Rommens. 2009. {``A Similarity Measure to Assess the Stability of
Classification Trees.''} \emph{Computational Statistics \& Data
Analysis} 53 (4): 1208--17.
\url{https://doi.org/10.1016/j.csda.2008.10.033}.

\bibitem[\citeproctext]{ref-delise2026era}
Delise, Timothy. 2026. {``Era Splitting -- Invariant Learning for
Decision Trees.''} \emph{Expert Systems with Applications} 299: 129687.
\url{https://doi.org/10.1016/j.eswa.2025.129687}.

\bibitem[\citeproctext]{ref-dietterich2000ensemble}
Dietterich, Thomas G. 2000. {``An Experimental Comparison of Three
Methods for Constructing Ensembles of Decision Trees: Bagging, Boosting,
and Randomization.''} \emph{Machine Learning} 40 (2): 139--57.
\url{https://doi.org/10.1023/A:1007607513941}.

\bibitem[\citeproctext]{ref-domingos2000vfdt}
Domingos, Pedro, and Geoff Hulten. 2000. {``Mining High-Speed Data
Streams.''} \emph{Proceedings of the Sixth ACM SIGKDD International
Conference on Knowledge Discovery and Data Mining (KDD '00)}, 71--80.
\url{https://doi.org/10.1145/347090.347107}.

\bibitem[\citeproctext]{ref-lee2024anls}
Lee, Joon, Min Kyu Sim, and Jin-Seon Hong. 2024. {``Assessing Decision
Tree Stability: A Comprehensive Method for Generating a Stable Decision
Tree.''} \emph{IEEE Access} 12: 90061--72.
\url{https://doi.org/10.1109/ACCESS.2024.3419228}.

\bibitem[\citeproctext]{ref-li2002instability}
Li, Rong-Hua, and Stewart Clarke. 2002. {``Instability of Decision Tree
Classification Algorithms.''} \emph{Proceedings of the Eighth ACM SIGKDD
International Conference on Knowledge Discovery and Data Mining (KDD
'02)}, 570--75. \url{https://doi.org/10.1145/775047.775131}.

\bibitem[\citeproctext]{ref-mirzamomen2017stable}
Mirzamomen, Zahra, and Mohammad Reza Kangavari. 2017. {``A Framework to
Induce More Stable Decision Trees for Pattern Classification.''}
\emph{Pattern Analysis and Applications} 20 (4): 991--1004.
\url{https://doi.org/10.1007/s10044-016-0542-2}.

\bibitem[\citeproctext]{ref-quinlan1986induction}
Quinlan, J. Ross. 1986. {``Induction of Decision Trees.''} \emph{Machine
Learning} 1 (1): 81--106. \url{https://doi.org/10.1023/A:1022643204877}.

\bibitem[\citeproctext]{ref-strobl2007bias}
Strobl, Carolin, Anne-Laure Boulesteix, and Thomas Augustin. 2007.
{``Unbiased Split Selection for Classification Trees Based on the {G}ini
Index.''} \emph{Computational Statistics \& Data Analysis} 52 (1):
483--501. \url{https://doi.org/10.1016/j.csda.2006.12.030}.

\bibitem[\citeproctext]{ref-turney1995stability}
Turney, Peter D. 1995. {``Bias and the Quantification of Stability.''}
\emph{Machine Learning} 20 (1): 23--33.
\url{https://doi.org/10.1007/BF00993473}.

\bibitem[\citeproctext]{ref-wang2017tsallis}
Wang, Yisen, Shu-Tao Xia, and Jia Wu. 2017. {``A Less-Greedy Two-Term
Tsallis Entropy Information Metric Approach for Decision Tree
Classification.''} \emph{Knowledge-Based Systems} 120: 34--42.
\url{https://doi.org/10.1016/j.knosys.2016.12.021}.

\bibitem[\citeproctext]{ref-yang2019pairgain}
Yang, Bo-Bo, Wei Gao, and Min Li. 2019. {``On the Robust Splitting
Criterion of Random Forest.''} \emph{2019 IEEE 19th International
Conference on Data Mining (ICDM)}, 1420--25.
\url{https://doi.org/10.1109/ICDM.2019.00184}.

\end{CSLReferences}

\end{document}
